\section{Introduction}\label{sec:introduction}

A quench starts a spin system from a simple nonequilibrium law and then
lets it evolve at a prescribed temperature. At the critical temperature
of the Sherrington--Kirkpatrick model, the relevant relaxation time grows
with the number of spins. We determine the two-time autocorrelation after
a quench from independent symmetric spins. The limit retains the random
spectral edge of the interaction matrix and is governed by the same
diffusion as the stationary critical autocorrelation.

\subsection{The model and the result}

Let \(D\) be either the standard Gaussian law or the uniform law on
\(\{-1,1\}\). For \(n\ge2\), let \(J_{ij}\), \(i<j\), be independent with
law \(D\), and put \(J_{ji}=J_{ij}\), \(J_{ii}=0\). The critical zero-field
Gibbs measure is
\begin{equation}\label{eq:gibbs}
 \pi_{n,J}(\sigma)
 =Z_{n,J}^{-1}
   \exp\left\{n^{-1/2}\sum_{i<j}J_{ij}\sigma_i\sigma_j\right\},
 \qquad \sigma\in\{-1,1\}^n .
\end{equation}
Every site has an independent rate-one clock. At a ring at site \(i\), its
spin is replaced by \(a\in\{-1,1\}\) with conditional probability
\[
 \frac{\exp(a h_i(\sigma))}{2\cosh h_i(\sigma)},
 \qquad
 h_i(\sigma)=n^{-1/2}\sum_{j\ne i}J_{ij}\sigma_j .
\]
Write \(P_t^J\) for this heat-bath semigroup. Thus one unit of time contains
\(n\) updates on average. Let \(\nu_n\) be the uniform law on the cube,
independent of \(J\), and set \(T_n=n^{2/3}\). The stationary and quench
autocorrelations in critical units are, respectively,
\begin{align}
 A_{n,J}(t)
   &=n^{-2/3}\sum_{i=1}^n
      \ip{\sigma_i}{P_{tT_n}^J\sigma_i}_{\pi_{n,J}},
       \qquad t>0,                                      \label{eq:finite-A}\\
 B_{n,J}(s,t)
   &=n^{-2/3}\sum_{i=1}^n
      \E_{\nu_n}^J[\sigma_i(sT_n)\sigma_i((s+t)T_n)],
       \qquad s,t>0.                                    \label{eq:finite-B}
\end{align}
Thermal and clock randomness are averaged in these functions; the
interaction \(J\) remains random. The variable \(s\) is the waiting time
after the quench and \(t\) is the subsequent lag.

We recall the edge data of the stationary theorem
\cite[Proposition~2.3, Proposition~5.1, and Theorem~6.8]{FUA}.
Let \(g=(g_a)_{a\ge1}\) be the increasing upper GOE edge process, the
joint limit in law of the leading coordinates: for each fixed \(m\),
\[
 \bigl(n^{2/3}(2-\lambda_a)\bigr)_{a=1}^m
       \ \weakto\ (g_a)_{a=1}^m ,
\]
where \(\lambda_1\ge\lambda_2\ge\cdots\) are the eigenvalues of a GOE
matrix with off-diagonal variance \(1/n\)
\cite[Theorem~1.1]{RamirezRiderVirag2011}. For any \(b+g_1>0\), put
\(\ell_a=(g_a+b)^{-1}\) and
\begin{equation}\label{eq:edge-D-intro}
 D_g(b)=\lim_{L\to\infty}
 \left\{\sqrt L-\sum_{a\ge1}
 \left(\frac1{g_a+b}-\frac1{g_a+L}\right)\right\}.
\end{equation}
The limit is in probability over \(g\), with the almost-sure version
chosen along a deterministic sequence as in the cited stationary theorem.
If \(Z_a\) are independent \(N(0,\ell_a)\), their centered square sum
converges. The continuous densities of this sum and its tails define
\[
 \Pi_g=\Law\left((Z_a)_{a\ge1}\,\middle|\,
              \sum_{a\ge1}(Z_a^2-\ell_a)=D_g(b)\right).
\]
This measure is independent of the admissible \(b\). The cylinder
gradient form
\[
 \cE_g(F)=\int\sum_{a\ge1}|\partial_aF|^2\,\dd\Pi_g
\]
is closable. Let \(H_g\) be the nonnegative operator of its closure, and
let \(c_*>0\) be the deterministic microscopic coefficient in
\cite[Theorem~6.8]{FUA}. We write
\begin{equation}\label{eq:edge-semigroup}
 S_t^g=e^{-t c_*H_g},
 \qquad
 A_g(t)=\sum_{a\ge1}\ip{x_a}{S_t^g x_a}_{\Pi_g},\quad t>0.
\end{equation}
The stationary theorem gives a Borel version of \(A_g\) in
\(C_{\mathrm{loc}}((0,\infty))\). Section~\ref{sec:background} records the
finite normalization and the stationary edge, static-coordinate, and
relaxed-energy inputs.

For an open set \(V\subset\R^d\), \(C_{\mathrm{loc}}(V)\) denotes the
space of continuous functions with uniform convergence on every compact
subset. Fix the countable family \(\mathscr V\) of bounded smooth cylinder
tests constructed in Section~\ref{sec:gaussian-limits}. It is dense in
\(L^2(\Pi_g)\) for each edge realization under consideration.
We call a family \(h_s^g\in L^2(\Pi_g)\) measurable
when \((g,s)\mapsto\Pi_g[h_s^g\phi]\) is measurable for each test
\(\phi\in\mathscr V\). These pairings determine each density class.

\newpage
\begin{theorem}[Critical quench universality]\label{thm:main}
There is a measurable family \((h_s^g)_{s>0}\), with
\(h_s^g\in L^2(\Pi_g)\), such that for almost every \(g\) the following
properties hold.
\begin{enumerate}
\item For every \(s>0\), \(h_s^g\ge0\) and \(\Pi_g[h_s^g]=1\). For
      \(s,t>0\),
      \[
       h_{s+t}^g=S_t^g h_s^g.
      \]
      Moreover,
      \[
       h_s^g\Pi_g\ \weakto\ \delta_0
       \quad\text{in the product topology on }\R^\N
       \quad(s\downarrow0),
       \qquad
       \|h_s^g-1\|_{L^1(\Pi_g)}\longrightarrow0
       \quad(s\to\infty).
      \]
\item The series
      \begin{equation}\label{eq:limit-B}
       B_g(s,t)=\sum_{a\ge1}
          \int h_s^g(x)\,x_a\,(S_t^g x_a)(x)\,\dd\Pi_g(x)
      \end{equation}
      converges locally uniformly on \((0,\infty)^2\) and defines a
      continuous function.
\item For both choices of \(D\),
      \[
       \Law_D(A_{n,J},B_{n,J})
       \ \weakto\
       \Law_g(A_g,B_g)
      \]
      in
      \(C_{\mathrm{loc}}((0,\infty))
        \times C_{\mathrm{loc}}((0,\infty)^2)\).
      The two coupling laws have the same \(g\), the same family
      \(h^g\), and the same coefficient \(c_*\).
\item For every \(0<a<b<\infty\),
      \[
       \sup_{a\le t\le b}|B_g(s,t)-A_g(t)|
          \longrightarrow0
       \quad\text{in probability over }g
       \quad(s\to\infty).
      \]
\end{enumerate}
The positive-time density limits of the Gaussian dynamics with initial
law \(\nu_n\) select this family, as made precise in
Theorem~\ref{thm:entrance-selection}. The normalization is exactly
\eqref{eq:finite-A}--\eqref{eq:finite-B}.
\end{theorem}

The boundary \(\delta_0\) is interpreted in the product topology.
At small times each fixed edge coordinate can vanish while the
renormalized constraint is carried by coordinates whose indices tend
to infinity. The proof identifies which such entrance is selected by
the uniform starts through a stronger, quantitative condition on these
coordinates.

\subsection{History and significance}

Sherrington and Kirkpatrick introduced the infinite-range spin-glass
model in 1975 \cite{SK1975}. The subsequent dynamical theory studied
memory through two-time correlation and response functions. In particular,
Sompolinsky and Zippelius developed a dynamical mean-field treatment
for soft spins \cite{SompolinskyZippelius1982}, and Cugliandolo and
Kurchan analyzed nonequilibrium relaxation after a quench in a
soft-spin Langevin formulation \cite{CugliandoloKurchan1994}. Their
asymptotic regime takes the system-size limit before the long-time limit
and differs from the finite-size critical window considered here.
For critical Ising heat-bath dynamics, Billoire and Campbell
\cite[Equation~(7)]{BilloireCampbell2011} numerically tested the
stationary finite-size form
\[
 q_c(n,t)\sim n^{-1/3}F(t/n^{2/3}),
\]
with time measured in updates per spin. Theorem~\ref{thm:main} studies
the two positive times following a uniform quench and identifies the
law of the sample-dependent limiting function.

The spectral mechanism is especially transparent in spherical spin
glasses. Kosterlitz, Thouless, and Jones
\cite{KosterlitzThoulessJones1976} analyzed the infinite-range Gaussian
spherical spin glass using properties of large random matrices.
Cugliandolo and Dean \cite{CugliandoloDean1995}
studied relaxation and aging for the spherical model, while
Ben Arous, Dembo, and Guionnet \cite{BenArousDemboGuionnet2001}
proved limiting dynamical and aging results with a soft spherical
constraint. Fyodorov, Perret, and Schehr
\cite{FyodorovPerretSchehr2015} later connected finite-size randomness
on the \(n^{2/3}\) scale to the edge in zero-temperature spherical
dynamics. These works concern different spin spaces or temperature
regimes; they explain why the critical edge should remain visible in
a nonequilibrium limit.

The stationary critical Ising theorem \cite{FUA} identifies the
conditioned edge measure, its closed dynamical form, and the
microscopic coefficient \(c_*\). Its equilibrium comparison uses
the conditional eigenframe averaging of Comets \cite[Equations~(2.1)--(2.2)]{Comets1996},
which equates the Haar-averaged cube and spherical partition integrals
when both use normalized reference measures and the same quadratic
Hamiltonian, and the
critical conditioned-Gaussian description of Du and Huang
\cite{DuHuang2026FreeEnergy,DuHuang2026Overlap}. Using the cap and
spectral estimates of \cite{FUA,CM,CS},
Sections~\ref{sec:edge-estimates} and~\ref{sec:warming} extend cap
control down to a tight lower threshold for \(np^3\) and derive warming
uniformly over the initial law at each positive critical time. The
order-\(n\) entropy argument uses Lemma~\ref{warm:posterior-mean}, which
strengthens the square-density posterior-mean estimate of
\cite[Proposition~7.3]{SG} to square-root dependence on the relative
heat-bath energy. The stationary relaxed-energy theorem is applied within its
polynomially bounded test class through bounded resolvent optimizers in
Section~\ref{sec:gaussian-limits}.

\subsection{The argument}

In the Gaussian model, adding an independent GOE diagonal changes the spin
Hamiltonian only by a constant. Let \(u_a\) be orthonormal eigenvectors of
the completed matrix and put \(x_a=n^{-1/3}\langle u_a,\sigma\rangle\).
If \(w_s\) is the Gibbs-relative
density at time \(sT_n\), orthogonality rewrites the quench correlation as
\[
 B_{n,J}(s,t)=\sum_{a=1}^n
       \pi_{n,J}[w_s\,x_a\,P_{tT_n}^Jx_a].
\]
This formula separates the law reached after the quench from subsequent
relaxation. The source identifies the relaxation operator for sufficiently
controlled tests. Three further arguments are needed.

We first regularize the law reached after the quench. At scale \(p\), a
nested Gaussian observation gives noisy information about the spin
projections onto modes with \(2-\lambda_a\le p^2\); conditioning caps
those eigenvalues at \(2-p^2\). On the retained spectral events, their
number is comparable to \(k_p=np^3\).
Section~\ref{sec:edge-estimates} bounds the variance left unresolved for
\(0\le f\le1\) by \(Cp^{-2}\) times the heat-bath energy plus
\(Ce^{-ck_p^a}\), down to a fixed threshold for \(k_p\) at each prescribed
disorder confidence. In the small-entropy part of
Section~\ref{sec:warming}, a contradiction argument assumes that a bounded
level test \(f\) has less energy than the claimed profile bound. At a later
scale \(R\), this makes the normalized tilt by \(f^2\) close to the
normalized observation tilt by \(F_R^2\), where
\(F_R=\E[f\mid h_R]\). Following its entropy through earlier observations,
together with the source path to a product posterior at higher entropy
levels, yields entropy and kernel bounds at every fixed positive critical
time, uniformly over the initial law.

For bounded cylinder sources, the finite Markov resolvent contraction
bounds the optimizers used in the stationary relaxed-energy theorem.
Section~\ref{sec:gaussian-limits} derives strong semigroup convergence
from that theorem and combines it with the warming bounds on densities
and transition kernels. This produces positive-time families on represented
Gaussian subsequences, evolving by \(S_t^g\), which may initially retain
randomness beyond \(g\).
To select one, choose a measurable admissible \(b\) with \(D_g(b)>0\),
put \(r_i=g_i+b\), and define
\[
 M_g(x)=\sum_i\frac{x_i^2}{2r_i}.
\]
We prove that \(0<M_g(x)<\infty\) for \(\Pi_g\)-almost every \(x\) and
therefore may define, on that full-measure set, the probability weights
\[
 p_i(x)=\frac{x_i^2}{2r_iM_g(x)}.
\]
For each represented Gaussian limit, the averaged uniform-start
estimate yields a deterministic sequence \(s_k\downarrow0\) along
which these weights leave every fixed finite set in probability under
\(h_{s_k}\Pi_g\), almost surely in the represented data.
Section~\ref{sec:entrance-selection} considers semigroup families whose
densities are locally uniformly bounded in \(L^\infty(\Pi_g)\) and satisfy
such sequential escape. Angular ratios and a common rescaling give linear
mode equations driven by independent Brownian motions with a common
diffusion factor. The square constraint gives a scalar Volterra equation
whose high-mode responses approach a common kernel. A stopped coupling
proves fixed-\(g\) uniqueness and the full \(s\downarrow0\) product-topology
boundary. The proof then constructs the common family measurably from \(g\).

Finally, we compare smooth functions of the stationary and quench traces
under entry replacement \cite{Chatterjee2006}. An order-\(k\) derivative
with respect to the normalized interaction \(J_{ij}/\sqrt n\) is
accompanied by \(n^{-k/2}\), while summing over entries contributes
order \(n^2\). The two coupling laws match through order three, so the
averaged fourth derivative must gain a vanishing factor.
Sections~\ref{sec:history-cutoffs}--\ref{sec:replacement} first impose
static overlap-tail estimates for every interaction matrix in a
cutoff's support. Under the associated normalized weighted path laws,
these estimates yield history-overlap bounds with failure probability
\(O(n^{-A})\) for any prescribed \(A>0\), after the cutoff parameters
are chosen. Forward and reverse conditional estimates for
endpoint-normalized likelihoods on a binary tree of transition intervals
then control the derivatives over the full critical time span. The
estimates also include derivatives of the cutoff weights. They
transfer the joint trace laws and the probability of retaining the
cutoffs to Rademacher disorder; smooth grid tests then transfer
tightness in the locally uniform topology.
Section~\ref{sec:completion} assembles the result.
